\section{Introduction}\label{sec:introduction}

The six-vertex model assigns an arrow to each edge of the square lattice,
subject to the rule that two arrows enter and two arrows leave each vertex.
Its height function is a discrete random surface. We prove the Gaussian
free field limit of this surface throughout the critical range
\(a=b=1\), \(0<c\le2\), for the plane state specified below.

\subsection{The plane law and the height convention}

The two vertices having both horizontal arrows directed inward, or both
directed outward, have weight \(c\); the other four vertices have weight
one. A finite configuration has the product of its vertex weights.
On
\[
 T_{M,L}=(\Z/M\Z)\times(\Z/L\Z),\qquad M,L\ge4,\quad L\ \text{even},
\]
condition every vertical column of horizontal edges to contain \(L/2\)
left arrows and \(L/2\) right arrows. This is a conserved sector of the row
transfer matrix. We first take \(M\to\infty\), and then even \(L\to\infty\),
in the local weak topology of arrow configurations. We denote the resulting
plane law by \(\P_c\). Existence in the previously untreated range is part
of the theorem.

Give the face on the left of an arrow height one more than the face on its
right. The ice rule makes this a well-defined height \(h\), modulo an
additive constant, on the faces of the plane. For \(x=(x_1,x_2)\in\R^2\),
let \(h(x)\) be the height of the face with bottom-left corner
\((\lfloor x_1\rfloor,\lfloor x_2\rfloor)\), and set
\[
 h_\delta(x)=h(x/\delta),\qquad \delta>0.
\]
This convention also fixes the value on grid lines.

Let \(\Gamma\) be the real plane Gaussian free field modulo constants,
normalized by the Green kernel
\begin{equation}\label{eq:green}
 G(x,y)=-\frac1{2\pi}\log|x-y|.
\end{equation}
Thus its covariance on smooth compactly supported tests of total integral
zero is the integral of \(G\) against their product.
For compactly supported finite signed measures \(\mu,\nu\) of total mass
zero, write
\[
 \mathcal E(\mu,\nu)=\iint G(x,y)\,\mu(\dd x)\nu(\dd y).
\]
In saying that \(\mu\) has finite energy, we require this to be an ordinary
finite signed-measure integral for \(\nu=\mu\). We do not use cancellation
between infinite positive and negative logarithmic energies.

For the negative regularity assertion we retain, in addition to the usual
distributional restriction spaces, the literal seminorm in
\cite[Definition~2.6(iii)]{DKLM2026}. Let \(\mathcal T_1(Q)\), with
\(Q=(-1,1)^2\), consist of compactly supported signed densities in \(Q\),
of integral zero, whose zero extensions are \(1\)-Lipschitz. Set
\begin{equation}\label{eq:source-holder}
 \|f\|_{\mathcal C^\alpha(U)}
 =\sup_{\substack{0<\varepsilon\le1,\ \varphi\in\mathcal T_1(Q)\\
                   \supp\varphi\subset U/\varepsilon}}
   \varepsilon^{-\alpha}
   \left|\int f(x/\varepsilon)\varphi(x)\,\dd x\right|.
\end{equation}
This is a seminorm of the global plane field: \(U\) indexes the allowed
tests, and the sampled points \(x/\varepsilon\) need not lie in \(U\).
It is distinct from the ordinary restriction topology on \(U\).
The space is the completion modulo the kernel of this seminorm.
The absolute value does not change the source convention because the test
family is symmetric. We prove convergence for this convention as well as
for the usual local restriction Hölder, Besov and Sobolev spaces.

\begin{theorem}\label{thm:main}
For every fixed \(c\in(0,2]\), the prescribed iterated balanced-torus limit
\(\P_c\) exists. Define
\begin{equation}\label{eq:multiplier}
 \Delta=1-\frac{c^2}{2},\qquad
 \sigma(c)^2=\frac2{\arccos\Delta}=\frac1{\arcsin(c/2)}.
\end{equation}
As \(\delta\downarrow0\) through arbitrary positive real values,
the fields \(h_\delta\) converge to
\(\sigma(c)\Gamma\) in each of the following senses.
\begin{enumerate}
\item For every finite family \(\mu_1,\ldots,\mu_m\) of compactly supported
finite signed measures of total mass zero and finite energy,
\[
 \bigl(\langle h_\delta,\mu_j\rangle\bigr)_{j=1}^m
 \ \Longrightarrow\
 \mathcal N\!\left(0,
       \bigl(\sigma(c)^2\mathcal E(\mu_i,\mu_j)\bigr)_{i,j=1}^m\right).
\]
\item For every bounded open \(U\subset\R^2\) and \(-1<\alpha<0\), there
is convergence in law, with tightness, in \(\mathcal C^\alpha(U)\) modulo
constants, and in the corresponding spaces
\[
 B^\alpha_{p,q}(U),\qquad W^{\alpha,p}(U),
 \qquad 1\le p<\infty,\quad 1\le q\le\infty.
\]
The usual local restriction Hölder topology is also included, and no
regularity of \(\partial U\) is required. One may fix a
smooth compactly supported \(\rho\) of integral one and replace each field
\(X\) by \(X-\langle X,\rho\rangle\), using the same pinning for the limit.
\item For every \(k\ge1\) and every tuple
\((u_i,u_i')_{i=1}^k\) with
\(\{u_i,u_i'\}\cap\{u_j,u_j'\}=\varnothing\) for \(i\ne j\),
\begin{equation}\label{eq:increment-main}
 \E_c\prod_{i=1}^k
    \bigl(h_\delta(u_i')-h_\delta(u_i)\bigr)
 \longrightarrow
 \sigma(c)^k\sum_{\pi}\prod_{\{i,j\}\in\pi}C_{ij},
\end{equation}
where the sum is over pairings of \(\{1,\ldots,k\}\), is zero for odd \(k\),
and
\[
 C_{ij}=G(u_i',u_j')-G(u_i',u_j)-G(u_i,u_j')+G(u_i,u_j).
\]
The convergence is uniform on compact sets of such tuples.
\end{enumerate}
\end{theorem}

The theorem concerns this particular zero-slope plane state. It makes no
uniqueness assertion among all Gibbs states or all slopes, and no estimate
uniform as \(c\downarrow0\).

\subsection{Prior work and the new ingredients}

The prediction that critical six-vertex height fluctuations have a Gaussian
free field limit connects this interacting lattice surface with conformal
universality. Earlier rigorous approaches include the dimer correspondence
at the free-fermion point \(c=\sqrt2\) and constructive renormalization in
a neighborhood of that point; see Kenyon~\cite{Kenyon2000},
Giuliani--Mastropietro--Toninelli~\cite{GMT2017}, and the account in
\cite[Section~2.4]{DKLM2026}.
Duminil-Copin, Kozlowski, Lammers and Manolescu proved the stated convergence
for \(\sqrt3\le c\le2\) and identified the multiplier
\eqref{eq:multiplier}~\cite[Theorem~2.8]{DKLM2026}.
Their work also provides the balanced plane state for \(1\le c\le2\),
its equilibrium characterization, and a conditional identification of the
multiplier from the curvature of the free energy
\cite[Theorem~2.2, the paragraph after (6), and Theorems~4.3--4.4]{DKLM2026}.
Theorem~\ref{thm:main} resolves the Gaussian free field scaling-limit
conjecture for this prescribed balanced plane law throughout \(0<c\le2\),
retaining these established parameter values.

The finite algebra in our argument is the algebraic Bethe ansatz.
We use scalar-product and norm determinants in their trigonometric and
diagonally twisted forms~\cite{Slavnov1989,KMT1999,KMST2005}.
Those finite identities are separate from the uniform thermodynamic
estimates proved here. We also use the local row algebra and folded angular
continuation mechanism in~\cite[Section on angular transfer]{AngularRepository2026}.
Its original parameter range does not cover the whole theorem.
Lemma~\ref{lem:scalar-continuation} checks the finite local continuation
identities for the present range, and Sections~\ref{sec:thermodynamic}
and~\ref{sec:angular} prove the limiting kernel estimates used here.

Three parts of the argument are useful beyond this particular organization
of the proof. First, a root-tail estimate and an edge maximum principle
give uniform inverses for Gaudin matrices and for mixed discrete--continuous
quadrature operators. Combined with finite-dimensional residue conditions,
these bounds control simultaneous escapes of several rapidities in
determinant formulas. Second, angular spectra split into two widely
separated scales. Bounds on all zeros outside the two scale windows
justify the operator limits that yield holomorphic and antiholomorphic
Gaussian currents. Third, the normalization for \(c<1\) follows from
matching a homogeneous stiffness lower bound with a staggered
vacuum-overlap upper bound. The latter turns an exact finite determinant
into a spectral packing inequality. Ordinary moment estimates and analytic
desmoothing then recover point-increment correlations and the full test
classes in Theorem~\ref{thm:main}.

\subsection{Proof overview}\label{sec:overview}

For the new range \(0<c<\sqrt3\), introduce
\[
 \lambda=\arccos(-\Delta)\in(\pi/3,\pi),
 \qquad \eta=\lambda/\pi.
\]
Section~\ref{sec:finite} defines commuting row operators and their diagonal
marks. Section~\ref{sec:thermodynamic} constructs their limiting matrix
elements, proves the necessary uniform tail bounds, and shows that words
with widely separated angular parameters decorrelate.

These matrix elements define a Hilbert space with commuting angular rows
in Section~\ref{sec:angular}. Their joint spectral values are inner functions
of a complex row parameter. A mark acts as a current. The finite mass
\(\kappa\) of the spectrum transverse to common angular dilation is the
coefficient of the logarithmic height covariance. Section~\ref{sec:normalization}
computes \(\kappa\) directly when \(c<1\); for \(c\ge1\) the cited conditional
normalization theorem applies once convergence has been established.

Section~\ref{sec:chiral} separates the angular spectrum into two scales and
proves the Wick rules for the limiting currents. Section~\ref{sec:probability}
passes from these currents to smooth tests, finite-energy measures,
negative regularity and off-diagonal point increments. The final section
assembles these conclusions with the established interval. Each
thermodynamic passage below specifies which size tends to infinity first;
none relies on exchanging the auxiliary limits.
